%% kapa: exam=1 round=32 year=2021 period=2 form=A subject=accounting subject_ko=회계학 seq=12 number=52
%% source: Q-Net 감정평가사 기출문제 (cst00309) artlSeq=5208397 file=extracted/1008/5208397/A형/2021년도 제32회 감정평가사 1차 2교시 A형.pdf
%% conversion: kordoc 4.24.51 -> tools/gen1.py (auto); review=proofread-sonnet; flags=none
\begin{question}{52}
투자부동산에 관한 설명으로 옳지 않은 것은?
\begin{choices}
  \item 소유 투자부동산은 최초 인식시점에 원가로 측정한다.
  \item 투자부동산을 후불조건으로 취득하는 경우의 원가는 취득시점의 현금가격상당액으로 한다.
  \item 투자부동산의 평가방법으로 공정가치모형을 선택한 경우, 감가상각을 수행하지 아니한다.
  \item 공정가치로 평가하게 될 자가건설 투자부동산의 건설이나 개발이 완료되면 해당일의 공정가치와 기존 장부금액의 차액은 기타포괄손익으로 인식한다.
  \item 재고자산을 공정가치로 평가하는 투자부동산으로 대체하는 경우, 재고자산의 장부금액과 대체시점의 공정가치의 차액은 당기손익으로 인식한다.
\end{choices}
\end{question}
\answer{4}
